% Compile with XeLaTeX.
\documentclass[a4paper,11pt]{article}
\usepackage[top=7.0mm,bottom=7.0mm,left=17.8mm,right=17.8mm]{geometry}
\usepackage{fontspec}
\setmainfont{Noto Serif CJK JP}
\setsansfont{Noto Serif CJK JP}
\setmonofont{Noto Serif CJK JP}
\XeTeXlinebreaklocale "ja"
\XeTeXlinebreakskip=0pt plus 1pt
\usepackage[unicode,hidelinks]{hyperref}
\hypersetup{pdftitle={応募状 — 東京大学医学部附属病院},pdfauthor={黄 万鴻}}
\pagestyle{empty}
\setlength{\parindent}{1em}
\setlength{\parskip}{1.8pt}
\setlength{\emergencystretch}{2em}
\widowpenalty=10000
\clubpenalty=10000
\newcommand{\letterdate}{2026年10月6日}

\begin{document}
\fontsize{10pt}{18.6pt}\selectfont
\noindent\hfill\letterdate\par
\noindent
東京大学医学部附属病院\\
医療運営課　患者支援チーム　御中

\vspace{-2pt}
\begin{center}
{\bfseries\fontsize{14.5pt}{17.5pt}\selectfont 応募状}
\end{center}
\vspace{-3pt}
\noindent\hfill 黄 万鴻\par
\vspace{0pt}

\noindent 拝啓

私は東京大学大学院情報理工学系研究科を修了後、国立精神・神経医療研究センター（NCNP）で計算論的神経科学の研究に従事いたしました。工学と自然科学を学び研究するなかで、現代の複雑な問題には、技術的解決だけでなく、制度、価値、正義、倫理、信頼、人と人との関係を含む理解が不可欠であると考えるようになりました。そのため退職後は独立して学習・研究を続け、この一年ほど、ガバナンス理論、正義論、倫理学、哲学、法学、政治経済学、社会学、人類学、国際関係論などを横断して研究してまいりました。

とりわけ私が問題意識を持ったのは、AIをはじめとする技術の急速な発展と記号・情報の膨張のなかで生じる、象徴の危機、信頼の危機、主体性の疎外などです。これらを個別の現象として扱うだけでは不十分であり、存在、認識、関係形成の基礎から捉え直す必要があると考え、生成的関係理論（Generative Relational Theory: GRT）を構築してきました。この基礎的な再構成は、各実践領域の理論的再構成へも展開しています。具体的には、Generative Relational Governance（GRG）、Generative Relational Jurisprudence（GRJ）、Generative Relational Peacekeeping、Generative Relational Economics（GR-Eco）、Generative Relational Historiography（GRH）、Generative Relational Epistemology（GRE）などを構築してきました。さらに、人道主義、援助理論、国際協力、平和構築などについても、GRTから派生する実践理論として研究を進めています。

また、独立して研究を続けるなかで、理論の構築だけでなく、その実践可能性を継続的に検討し、研究成果の公開や将来の社会実践につなげるための基盤づくりも個人で進めてきました。私は、理論は実践から切り離されるべきではなく、具体的な人、制度、制約、責任に向き合うなかで実践知として鍛えられるべきだと考えています。ボランティア活動にも深い敬意を持っており、私自身、国連ボランティア（UNV）を通じて、SDG 3「すべての人に健康と福祉を」に関わるボランティア活動に参加した経験があります。

本職を志望するのは、これまでの研究・実践活動で培った情報整理、文書作成、企画、調整、IT活用の力を、患者支援とボランティア活動を支える具体的な仕事に結びつけたいからです。まず一人の職員として、患者様、ボランティア活動員、職員の皆様に対する責任を誠実に果たしたいと考えております。そのうえで、これまで形成してきた理論を実践のなかで用い、実践を通じてその妥当性と限界を検討し、得られた実践知によって理論を修正し、さらに発展させていきたいと考えております。

中国古典には、「上善は水の若し。水は善く万物を利して争わず、故に道に幾し」という言葉があります。また、「生じて有せず、為して恃まず、長じて宰せず、これを玄徳という」とも説かれています。私は、他者を利しながら争わず、何かを生み出しても自らの所有とせず、人を支えながら支配しないというこれらの言葉が、ボランティアという営みの精神と深く響き合うように感じています。そのような営みを担い、また支えておられる方々に、心より敬意を表します。

深秋の十月、貴院の皆様、そして世界のあらゆる場所で、人々の福祉と世界のよりよい未来のために力を尽くしておられる一人ひとりが、その献身とともに、自らの幸福をも得られることを心より願っております。

\noindent\hfill 敬具

\vspace{1pt}
\begin{flushright}
黄 万鴻\\
\href{mailto:huangwanhong@serendip.ngo}{huangwanhong@serendip.ngo}\\
\href{https://huang-wanhong.com}{huang-wanhong.com}
\end{flushright}
\end{document}
